\begin{frame}{Brecha tributaria}
\begin{itemize}
    \item Brecha tributaria = brecha de cumplimiento + brecha de política
    \item Brecha de cumplimiento (evasión) = 8,6\% del PIB.
    \item Brecha de política (gasto tributario) = 8,5\% del PIB.
    \item \textbf{Bajo el marco tributario de referencia, el GNC debería estar recaudando más del 30\% del PIB}
\end{itemize}
\end{frame}

\begin{frame}
\begin{table}[htbp]
\centering
\caption{Evasión por tipo de impuesto}
\label{tab:evasion_recaudo}
\large
\setlength{\tabcolsep}{7pt}
\renewcommand{\arraystretch}{1.22}
\begin{tabular}{lccc}
\toprule
Impuesto
& \shortstack{Recaudo\\(2025)\\\% del PIB}
& \shortstack{Evasión\\(2019)\\\% del PIB}
& \shortstack{Evasión\\(2019)\\\% del recaudo$^{*}$} \\
\midrule
IVA & 5,3 & \textbf{1,8} & 34\% \\
\shortstack[l]{Renta de personas\\jurídicas} & 6,0 & \textbf{5,7} & 95\% \\
\shortstack[l]{Renta de personas\\naturales} & 1,5 & \textbf{1,1} & 73\% \\
\bottomrule
\end{tabular}
\vspace{0.35cm}

\par\scriptsize{$^{*}$ Cociente entre la evasión estimada para 2019 y el recaudo de 2025. El porcentaje muestra un orden de magnitud y no una comparación para un mismo año.}

\vspace{0.12cm}
\par\scriptsize{\textit{Fuentes: Evasión con base en Lora et al. (2021). Recaudo con cálculos propios a partir del Balance del GNC del Ministerio de Hacienda.}}
\end{table}

\note{[Sources] Evasión: Lora et al. (2021), estimaciones para 2019. Recaudo: cálculos propios a partir del Balance del GNC del Ministerio de Hacienda, cifras para 2025. Los porcentajes del recaudo se calculan como evasión dividida por el recaudo de 2025.}

\end{frame}

\begin{frame}
\begin{table}[htbp]
\centering
\caption{Gasto tributario por tipo de impuesto}
\label{tab:gasto_tributario_recaudo}
\large
\setlength{\tabcolsep}{7pt}
\renewcommand{\arraystretch}{1.22}
\begin{tabular}{lccc}
\toprule
Impuesto
& \shortstack{Recaudo\\(2025)\\\% del PIB}
& \shortstack{Gasto tributario\\(2023)\\\% del PIB}
& \shortstack{Gasto tributario\\(2023)\\\% del recaudo$^{*}$} \\
\midrule
IVA & 5,3 & \textbf{5,6} & 106\% \\
\shortstack[l]{Renta de personas\\jurídicas} & 6,0 & \textbf{1,3} & 22\% \\
\shortstack[l]{Renta de personas\\naturales} & 1,5 & \textbf{1,6} & 107\% \\
\bottomrule
\end{tabular}
\vspace{0.35cm}

\par\scriptsize{$^{*}$ Cociente entre el gasto tributario estimado para 2023 y el recaudo de 2025. El porcentaje muestra un orden de magnitud y no una comparación para un mismo año.}

\vspace{0.12cm}
\par\scriptsize{\textit{Fuentes: Gasto tributario con base en el MFMP 2025. Recaudo con cálculos propios a partir del Balance del GNC del Ministerio de Hacienda.}}
\end{table}

\note{[Sources] Gasto tributario: Marco Fiscal de Mediano Plazo 2025, cifras para 2023. Recaudo: cálculos propios a partir del Balance del GNC del Ministerio de Hacienda, cifras para 2025. Los porcentajes del recaudo se calculan como gasto tributario dividido por el recaudo de 2025.}

\end{frame}

\begin{frame}{Brecha tributaria}
    \begin{itemize}
        \item Cerrar la brecha de cumplimiento: Problema administrativo (¿y político?)
        \item  Cerrar la brecha de política: Problema legal y jurídico:
        \begin{itemize}
            \item Sobretasas.
            \item Marchitar los beneficios tributarios que no esten: 
            \begin{itemize}
                \item Especificados en el Estatuto Tributario.
                \item Cuantificados en el PGN.
            \end{itemize}
        \end{itemize}
        \item Articulación: El marchitamiento de beneficios simplifica la tributación y facilita la reducción de tarifas, lo cual ayuda con el cumplimiento.
        \item El marchitamiento es, sin embargo, un desafío jurídico y político enorme.
    \end{itemize}
\end{frame}
